% SHARD-ID: AQM-44-TRACE-GAUGES-FULL-CONTINUUM-TOWERS
% SHARD-TITLE: Trace Gauges for Full Continuum Towers
% SHARD-SUMMARY: Defines coherent trace-density systems for finite-field scale towers.
% SHARD-SUMMARY: Proves that coherent trace data always exist and give coherent Weyl-net embeddings.
% SHARD-SUMMARY: Identifies the arbitrariness as trace gauge freedom and records the Frobenius obstruction.
% SHARD-KEYWORDS: continuum limit, trace gauge, finite field, Weyl net, Frobenius, coherence

\section{Trace Gauges for Full Continuum Towers}
\label{sec:trace-gauges-full-continuum-towers}

\subsection{Status and Source Anchors}

AQM-43 proved that pairwise trace-one elements repair the Weyl phase
obstruction for finite-field scale embeddings.  This shard proves that the
pairwise choices can be made coherently for the whole divisibility tower, and
records exactly what is gauge-like about the arbitrariness.  The source
anchors are Stacks Algebra,
\path{references/algebraic_geometry/stacks_project_algebra.tex}, lines
45278--45284 for nondegenerate separable trace pairing, and Stacks Varieties,
\path{references/algebraic_geometry/stacks_project_varieties.tex}, lines
7005--7112 for Frobenius.  All coherence statements below are local
derivations from those facts.

\subsection{Lamport Dependency Skeleton}

\[
\begin{array}{rcl}
D1&:&k_r=\mathbb F_{q^r}\subset\bar k,\quad r\mid s\Rightarrow k_r\subset k_s.\\
D2&:&\text{A trace density is }a_r\in k_r\text{ with }
      \Tr_{k_s/k_r}(a_s)=a_r.\\
L1&:&\Tr_{k_s/k_r}:k_s\to k_r\text{ is surjective.}\\
T1&:&\text{A coherent trace-density system }(a_r)_r\text{ with }a_1=1
      \text{ exists.}\\
D3&:&c_{s,r}=a_s/a_r.\\
L2&:&\Tr_{k_s/k_r}(c_{s,r})=1,\quad
      c_{t,r}=c_{t,s}c_{s,r}.\\
T2&:&(c_{s,r})\text{ gives coherent full-tower Weyl-net embeddings.}\\
T3&:&\text{The choice of }(a_r)\text{ is trace gauge data.}\\
T4&:&p\mid s/r\text{ forbids a relative-Frobenius-fixed trace-one scalar.}
\end{array}
\]

\subsection{Coherent Trace Densities D1--T1}

Fix \(k=\mathbb F_q\) and \(k_r=\mathbb F_{q^r}\subset\bar k\).  A coherent
trace-density system is a family \(a_r\in k_r\) such that \(a_1=1\) and,
whenever \(r\mid s\),
\[
  \Tr_{k_s/k_r}(a_s)=a_r .
\]
This condition forces \(a_r\ne0\), because
\[
  \Tr_{k_r/k}(a_r)=a_1=1.
\]

\paragraph{Lemma \(L1\).}
For \(r\mid s\), the trace map
\[
  \Tr_{k_s/k_r}:k_s\to k_r
\]
is surjective.

\emph{Proof.}
Finite fields are perfect, hence \(k_s/k_r\) is separable.  The sourced
separable trace pairing is nondegenerate:
\[
  (x,y)\longmapsto \Tr_{k_s/k_r}(xy).
\]
If the trace map itself were zero, then \(\Tr_{k_s/k_r}(1\cdot y)=0\) for
every \(y\), contradicting nondegeneracy with \(x=1\).  Thus the trace is a
nonzero \(k_r\)-linear map from \(k_s\) to the one-dimensional \(k_r\)-space
\(k_r\), so it is surjective.

\paragraph{Theorem \(T1\).}
Coherent trace-density systems exist.

\emph{Proof.}
Let \(N_m=m!\).  Start with \(a_{N_1}=a_1=1\).  Having chosen
\(a_{N_m}\), Lemma \(L1\) gives \(a_{N_{m+1}}\in k_{N_{m+1}}\) with
\[
  \Tr_{k_{N_{m+1}}/k_{N_m}}(a_{N_{m+1}})=a_{N_m}.
\]
For an arbitrary \(r\ge1\), choose \(m\) with \(r\mid N_m\), and define
\[
  a_r=\Tr_{k_{N_m}/k_r}(a_{N_m}).
\]
This is independent of \(m\): if \(n\ge m\), then transitivity of trace and
the recursive construction give
\[
  \Tr_{k_{N_n}/k_r}(a_{N_n})
  =
  \Tr_{k_{N_m}/k_r}
  \bigl(\Tr_{k_{N_n}/k_{N_m}}(a_{N_n})\bigr)
  =
  \Tr_{k_{N_m}/k_r}(a_{N_m}).
\]
If \(r\mid s\), choose \(m\) with \(s\mid N_m\).  Then
\[
  \Tr_{k_s/k_r}(a_s)
  =
  \Tr_{k_s/k_r}
  \bigl(\Tr_{k_{N_m}/k_s}(a_{N_m})\bigr)
  =
  \Tr_{k_{N_m}/k_r}(a_{N_m})
  =
  a_r.
\]
Thus \((a_r)_r\) is coherent.

\subsection{Coherent Weyl Embeddings D3--T2}

Given a coherent trace-density system, define
\[
  c_{s,r}=a_s/a_r\in k_s
\]
for \(r\mid s\), where \(a_r\) is embedded in \(k_s\).

\paragraph{Lemma \(L2\).}
\[
  \Tr_{k_s/k_r}(c_{s,r})=1,\qquad
  c_{t,r}=c_{t,s}c_{s,r}\quad(r\mid s\mid t).
\]

\emph{Proof.}
Since \(a_r\in k_r^\times\),
\[
  \Tr_{k_s/k_r}(a_s/a_r)
  =
  a_r^{-1}\Tr_{k_s/k_r}(a_s)=1.
\]
The multiplicative identity is immediate:
\[
  c_{t,s}c_{s,r}=(a_t/a_s)(a_s/a_r)=a_t/a_r=c_{t,r}.
\]

Let \(X_0=\Spec R\) be affine over \(k\), and put
\(X_r=X_0\times_k k_r\).  If \(x\in |X_r|_{\rm cl}\) has residue field
\(K=\kappa(x)\), then the points of \(X_s\) above \(x\) are encoded by
\[
  B_{s,r,x}=K\otimes_{k_r}k_s .
\]
Use \(1\otimes c_{s,r}\in B_{s,r,x}\) as the trace-normalizing scalar.  Its
trace to \(K\) is \(1\): choose a \(k_r\)-basis of \(k_s\); the multiplication
matrix of \(c_{s,r}\) has trace \(\Tr_{k_s/k_r}(c_{s,r})=1\), and base change
of that matrix from \(k_r\) to \(K\) leaves the trace equal to \(1\).

\paragraph{Theorem \(T2\).}
The coherent trace-density system gives coherent monomorphisms of all
open-local Weyl nets
\[
  \alpha_{s,r,U}:\mathcal A_r(U)\hookrightarrow
  \mathcal A_s(\pi_{s,r}^{-1}U),
  \qquad r\mid s.
\]

\emph{Proof.}
AQM-43 proves that any trace-one normalizing scalar gives a Weyl-compatible
monomorphism at one transition.  The preceding paragraph gives trace-one
scalars at every closed point, and Lemma \(L2\) gives
\[
  (1\otimes c_{t,s})(1\otimes c_{s,r})=1\otimes c_{t,r}.
\]
Therefore the label embeddings compose strictly, hence the induced Weyl-net
monomorphisms compose strictly.

\subsection{Gauge Meaning and Frobenius T3--T4}

The family \((a_r)\), equivalently \((c_{s,r})\), is called a trace gauge.  It
is not extra geometry of \(X_0\); it is extra presentation data needed to
compare Weyl phases across scales.  The prime-to-\(p\) subsystem has the
canonical gauge \(c_{s,r}=(s/r)^{-1}\).  The full tower has many gauges, by
Theorem \(T1\), but no canonical one is singled out by the finite-field
embeddings alone.

Frobenius acts on gauges.  If \(c_{s,r}\) has trace \(1\), then
\[
  \Tr_{k_s/k_r}(c_{s,r}^q)=\Tr_{k_s/k_r}(c_{s,r})^q=1.
\]
Thus Frobenius sends one trace gauge to another.  For \(p\mid s/r\), however,
there is no trace-one scalar fixed by the relative Frobenius group.  If
\(c\in k_s\) is fixed by \(\operatorname{Gal}(k_s/k_r)\), then \(c\in k_r\),
and
\[
  \Tr_{k_s/k_r}(c)=(s/r)c=0
\]
in characteristic \(p\).  This cannot equal \(1\).

For \(k=\mathbb F_3\), the first such case is
\(\mathbb F_3\subset\mathbb F_{27}\).  A trace-one \(c\) exists, but it cannot
lie in \(\mathbb F_3\).  Hence the full tower supports Frobenius only
gauge-covariantly unless an additional Frobenius-compatible gauge convention
is supplied.  This is the precise sense in which the arbitrariness is gauge
freedom rather than a mathematical obstruction to the continuum limit.
